\documentclass[../../main.tex]{subfiles}

\begin{document}

\section{累积分布函数与 Borel 测度}\label{section:ch20-累积分布}

令 $I=[a,b]$ 为实数的闭有界区间而 $\mathcal{B}(I)$ 是 $I$ 的 Borel 子集族. 我们称 $\mathcal{B}(I)$ 上的有限测度 $\mu$ 为 Borel 测度. 对这样的测度, 定义函数 $g_{\mu}:I\to\mathbf{R}$ 为
\begin{align*}
\text{对所有 }I\text{ 中的 }x,\quad g_{\mu}(x)=\mu[a,x].
\end{align*}
函数 $g_{\mu}$ 称为 $\mu$ 的累积分布函数.

\begin{proposition}\label{proposition:ch20-累积分布对应}
令 $\mu$ 为 $\mathcal{B}(I)$ 上的 Borel 测度. 则它的累积分布函数 $g_{\mu}$ 是递增与右连续的. 反过来, 每个递增与右连续的函数 $g:I\to\mathbf{R}$ 是 $\mathcal{B}(I)$ 上唯一的 Borel 测度 $\mu_{g}$ 的累积分布函数.
\end{proposition}

\begin{proof}
首先令 $\mu$ 为 $\mathcal{B}(I)$ 上的 Borel 测度. 它的累积分布函数当然是递增与有界的. 令 $x_{0}$ 属于 $[a,b)$ 而 $\{x_{k}\}$ 是属于 $(x_{0},b]$ 且收敛于 $x_{0}$ 的递减序列. 则 $\bigcap_{k=1}^{\infty}(x_{0},x_{k}]=\varnothing$. 因此由于 $\mu$ 是有限的, 根据测度的连续性,
\begin{align*}
0=\mu(\varnothing)=\lim_{k\to\infty}\mu(x_{0},x_{k}]=\lim_{k\to\infty}[g_{\mu}(x_{k})-g_{\mu}(x_{0})].
\end{align*}
因此 $g_{\mu}$ 在 $x_{0}$ 右连续.

为证明反命题, 令 $g:I\to\mathbf{R}$ 为递增的右连续函数. 考虑空集、单点集 $\{a\}$ 以及所有 $I$ 的所有形如 $(c,d]$ 的子区间组成的子集族 $\mathcal{S}$. 则 $\mathcal{S}$ 是半环. 考虑集函数 $\mu:\mathcal{S}\to\mathbf{R}$, 它定义为 $\mu(\varnothing)=0$, $\mu\{a\}=g(a)$ 且
\begin{align*}
\text{对 }(c,d]\subseteq I,\quad\mu(c,d]=g(d)-g(c).
\end{align*}
我们将以下结论的证明留作练习: 若 $(c,d]\subseteq I$ 是有限不交族 $\bigcup_{k=1}^{n}(c_{k},d_{k}]$ 的并, 则
\begin{align*}
g(d)-g(c)=\sum_{k=1}^{n}[g(d_{k})-g(c_{k})].
\end{align*}
且若 $(c,d]\subseteq I$ 被可数族 $\bigcup_{k=1}^{\infty}(c_{k},d_{k}]$ 覆盖, 则
\begin{align}\label{align:ch20-预测度不等式}
g(d)-g(c)\leqslant\sum_{k=1}^{\infty}[g(d_{k})-g(c_{k})].
\end{align}
这意味着 $\mu$ 是半环 $\mathcal{S}$ 上的预测度. 根据 Carath\'eodory-Hahn 定理, 由 $\mu$ 诱导的 Carath\'eodory 测度 $\overline{\mu}$ 是 $\mu$ 的延拓. 特别地, $[a,b]$ 的每个开子集是 $\mu^{*}$ 可测的. 根据 Borel $\sigma$ 代数的最小性, $\mu^{*}$ 可测集的 $\sigma$ 代数包含 $\mathcal{B}(I)$. 函数 $g$ 是 $\overline{\mu}$ 在 $\mathcal{B}(I)$ 上的限制的累积分布函数. 由于对每个 $x\in[a,b]$
\begin{align*}
\mu[a,x]=\mu\{a\}+\mu(a,x]=g(a)+[g(x)-g(a)]=g(x).
\end{align*}
\end{proof}

将 Borel 测度的连续性与它的累积分布函数的连续性联系在一起是自然的. 我们有以下令人满意的关系, 其证明留作练习.

\begin{proposition}\label{proposition:ch20-绝对连续}
令 $\mu$ 为 $\mathcal{B}(I)$ 上的 Borel 测度而 $g_{\mu}$ 是它的累积分布函数. 则测度 $\mu$ 关于 Lebesgue 测度绝对连续当且仅当函数 $g_{\mu}$ 是绝对连续的.
\end{proposition}

对于 $[a,b]$ 上的有界 Lebesgue 可测函数 $f$, Lebesgue 积分有定义, 其中 $m$ 表示 Lebesgue 测度. 对于 $[a,b]$ 上的有界函数 $f$, 其不连续点的集合的 Lebesgue 测度为零, 我们证明 Riemann 积分 $\int_{a}^{b}f(x)dx$ 有定义且
\begin{align*}
\int_{[a,b]}f\,dm=\int_{a}^{b}f(x)\,dx.
\end{align*}
这些积分有两个推广--Lebesgue-Stieltjes 积分与 Riemann-Stieltjes 积分, 我们现在简要地讨论它们.

令函数 $g:I\to\mathbf{R}$ 为递增与右连续的. 对有界 Borel 可测函数 $f:I\to\mathbf{R}$, 我们定义 $f$ 在 $[a,b]$ 上关于 $g$ 的 Lebesgue-Stieltjes 积分 $\int_{[a,b]}f\,dg$ 为
\begin{align}\label{align:ch20-lebesgue-stieltjes}
\int_{[a,b]}f\,dg=\int_{[a,b]}f\,d\mu_{g}.
\end{align}
现在假定 $f$ 是有界 Borel 可测函数而 $g$ 是递增与绝对连续的. 则 $g'$ 是 $[a,b]$ 上的 Lebesgue 可积函数, 而 $fg'$ 也如此. 我们有
\begin{align}\label{align:ch20-绝对连续stieltjes}
\int_{[a,b]}f\,dg=\int_{[a,b]}fg'\,dm.
\end{align}
其中右边的积分是 $fg'$ 关于 Lebesgue 测度 $m$ 的积分. 为证明这个公式, 观察到它对 $f$ 是 Borel 简单函数成立, 接着根据简单逼近定理与 Lebesgue 控制收敛定理, 它对有界 Borel 可测函数 $f$ 成立. 在这种情形, 由命题\cref{proposition:ch20-绝对连续}, $\mu_{g}$ 关于 $m$ 是绝对连续的. 我们把证明函数 $g'$ 是 $\mu_{g}$ 的关于 $m$ 的 Radon-Nikodym 导数留给读者(见习题 44).

Riemann-Stieltjes 积分以与 Lebesgue-Stieltjes 积分推广 Lebesgue 积分相同的方式推广 Riemann 积分. 我们简要地描述这个推广. 若 $P=\{x_{0},x_{1},\cdots,x_{n}\}$ 是 $[a,b]$ 的分划, 我们令 $\|P\|$ 表示由 $P$ 确定的区间的最大长度, 而令 $C=\{c_{1},\cdots,c_{n}\}$, 其中每个 $c_{i}$ 属于 $[x_{i-1},x_{i}]$. 对两个有界函数 $f:[a,b]\to\mathbf{R}$ 与 $g:[a,b]\to\mathbf{R}$, 考虑形如
\begin{align*}
S(f,g,P,C)=\sum_{k=1}^{n}f(c_{i})\cdot[g(x_{i})-g(x_{i-1})]
\end{align*}
的和式. 若存在实数 $A$ 使得对每个 $\varepsilon>0$, 存在 $\delta>0$ 使得
\begin{align*}
\text{若 }\|P\|<\delta,\quad\text{则 }|S(f,g,P,C)-A|<\varepsilon,
\end{align*}
则可以认为 $f$ 在 $I$ 上关于 $g$ Riemann-Stieltjes 可积且设
\begin{align*}
A=\int_{a}^{b}f(x)\,dg(x).
\end{align*}
显然, 若对所有 $x\in[a,b]$, $g(x)=x$, 则 $f$ 关于 $g$ 的 Riemann-Stieltjes 积分就是 $f$ 的 Riemann 积分. 此外, 若 $f$ 是连续的而 $g$ 是单调的, 则 $f$ 在 $I$ 上关于 $g$ 是 Riemann-Stieltjes 可积的. Camille-Jordan 的定理告诉我们有界变差函数是递增函数的差. 因此 $I$ 上的连续函数关于有界变差函数在 $I$ 上是 Riemann-Stieltjes 可积的. Lebesgue-Stieltjes 积分与 Riemann-Stieltjes 积分对不同的函数类定义. 然而, 若 $f$ 连续而 $g$ 递增且绝对连续, 这两种积分都有定义, 在这种情形它们是相等的, 即
\begin{align*}
\int_{a}^{b}f(x)\,dg(x)=\int_{[a,b]}f\,dg.
\end{align*}
因为(见习题 36 与 37)每个积分等于 $\int_{[a,b]}fg'\,dm$, 即 $fg'$ 在 $[a,b]$ 上关于 Lebesgue 测度 $m$ 的 Lebesgue 积分.

\subsection*{习题}

\begin{exercise}\label{exercise:ch20-35}
证明命题\cref{proposition:ch20-绝对连续}.
\end{exercise}

\begin{exercise}\label{exercise:ch20-36}
假定 $f$ 是 $[a,b]$ 上的有界 Borel 可测函数, 而 $g$ 在 $[a,b]$ 上是递增与绝对连续的. 证明: 若 $m$ 表示 Lebesgue 测度, 则
\begin{align*}
\int_{[a,b]}f\,dg=\int_{[a,b]}fg'\,dm.
\end{align*}
\end{exercise}

\begin{exercise}\label{exercise:ch20-37}
假定 $f$ 是 $[a,b]$ 上的连续函数, $g$ 在 $[a,b]$ 上是递增与绝对连续的. 证明: 若 $m$ 表示 Lebesgue 测度, 则
\begin{align*}
\int_{a}^{b}f(x)\,dg(x)=\int_{[a,b]}fg'\,dm.
\end{align*}
\end{exercise}

\begin{exercise}\label{exercise:ch20-38}
令 $f$ 和 $g$ 为 $[-1,1]$ 上的函数使得在 $[-1,0]$ 上 $f=0$, 在 $(0,1)$ 上 $f=1$. 证明 $f$ 关于 $g$ 在 $[-1,1]$ 上不是 Riemann-Stieltjes 可积的, 但在 $[-1,0]$ 与 $[0,1]$ 上关于 $g$ 是 Riemann-Stieltjes 可积的.
\end{exercise}

\begin{exercise}\label{exercise:ch20-39}
证明不等式\eqref{align:ch20-预测度不等式}. (提示: 选取 $\varepsilon>0$. 根据 $g$ 的右连续性, 选取 $\eta_{i}>0$ 使得 $g(b_{i}+\eta_{i})<g(b_{i})+\varepsilon 2^{-i}$, 且选取 $\delta>0$ 使得 $g(a+\delta)<g(a)+\varepsilon$. 则开区间 $(a_{i},b_{i}+\eta_{i})$ 覆盖闭区间 $[a+\delta,b]$.)
\end{exercise}

\begin{exercise}\label{exercise:ch20-40}
对递增函数 $g:[a,b]\to\mathbf{R}$, 定义
\begin{align*}
g^{*}(x)=\lim_{y\to x^{+}}g(y).
\end{align*}
证明 $g^{*}$ 是右连续的递增函数且在 $g$ 的右连续点与 $g$ 一致. 除可数个点外, $g=g^{*}$. 证明 $(g^{*})^{*}=g^{*}$, 以及若 $g$ 和 $G$ 是递增函数且在使得它们都连续的点相等, 则 $g^{*}=G^{*}$. 若 $f$ 是 $[a,b]$ 上的有界 Borel 可测函数, 证明
\begin{align*}
\int_{[a,b]}f\,dg=\int_{[a,b]}f\,dg^{*}.
\end{align*}
\end{exercise}

\begin{exercise}\label{exercise:ch20-41}
(i) 证明每个有界变差的有界函数 $g$ 对应有限带号 Borel 测度 $\nu$, 使得
\begin{align*}
\text{对所有 }(c,d]\subseteq[a,b],\quad\nu(c,d]=g(d^{+})-g(c^{+}).
\end{align*}
(ii) 将 Lebesgue-Stieltjes 积分 $\int_{[a,b]}f\,dg$ 的定义推广到有界变差函数 $g$ 与有界 Borel 可测函数 $f$.
(iii) 证明: 若在 $[a,b]$ 上 $|f|\leqslant M$ 且若 $g$ 的全变差是 $T$, 则 $|\int_{[a,b]}f\,dg|\leqslant MT$.
\end{exercise}

\begin{exercise}\label{exercise:ch20-42}
令 $g$ 为 $[a,b]$ 上的递增函数, 满足 $g(a)=c$, $g(b)=d$, 且令 $f$ 为 $[c,d]$ 上的非负 Borel 可测函数. 证明
\begin{align*}
\int_{[a,b]}f\circ g\,dg=\int_{[c,d]}f\,dm.
\end{align*}
\end{exercise}

\begin{exercise}\label{exercise:ch20-43}
令 $g$ 为 $[a,b]$ 上的递增函数. 找到 $\mathcal{B}([a,b])$ 上的 Borel 测度 $\mu$ 使得
\begin{align*}
\text{对所有 }f\in C[a,b],\quad\int_{a}^{b}f(x)\,dg(x)=\int_{[a,b]}f\,d\mu.
\end{align*}
\end{exercise}

\begin{exercise}\label{exercise:ch20-44}
若 Borel 测度 $\mu$ 关于 Lebesgue 测度绝对连续, 证明它的 Radon-Nikodym 导数是它的累积分布函数的导数.
\end{exercise}

\begin{exercise}\label{exercise:ch20-45}
对 $\mathbf{R}$ 的所有 Borel 子集所成的族 $\mathcal{B}(\mathbf{R})$ 上的有限测度 $\mu$, 定义 $g:\mathbf{R}\to\mathbf{R}$ 为 $g(x)=\mu(-\infty,x]$. 证明每个满足 $\lim_{x\to-\infty}g(x)=0$ 的有界递增右连续函数 $g:\mathbf{R}\to\mathbf{R}$ 是 $\mathcal{B}(\mathbf{R})$ 上唯一的有限 Borel 测度的累积分布函数.
\end{exercise}

\end{document}